\section*{PLANTILLA DE CAPTURA / ELICITACIÓN\\DE REQUERIMIENTOS}
\begin{center}
\textit{Fase de Descubrimiento y Levantamiento de Información}
\end{center}

\section{Información de la Sesión / Entrevista}

\begin{tabularx}{\textwidth}{>{\bfseries}p{0.30\textwidth}X}
\toprule
\rowcolor{azulprincipal}\textcolor{white}{Campo} & \textcolor{white}{Detalle / Registro del Analista}\\
\midrule
Fecha de la sesión & 31 de julio de 2026\\
Analista líder & Juan Fernando Sique Muralles\\
Stakeholders / clientes & Elena Ramos, Gerente General\\
Proyecto / sistema & Sistema de Gestión de Citas Médicas\\
Objetivo de la reunión & Detallar las solicitudes requeridas para el sistema de gestión de citas médicas.\\
Código / referencia de sesión & SES-001\\
\bottomrule
\end{tabularx}

\section{Transcripción Cruda y Necesidades del Usuario}

\textit{Instrucciones: Anote en esta sección las frases textuales, dolores, flujos actuales o deseos expresados por el usuario, sin importar que contengan ambigüedades técnicas.}

\begin{longtable}{>{\bfseries}p{0.24\textwidth}p{0.64\textwidth}}
\toprule
\rowcolor{azulprincipal}\textcolor{white}{Ref. / Stakeholder} & \textcolor{white}{Necesidad / Dolor Expresado (Lenguaje Coloquial)}\\
\midrule
GER-01 (Gerente General) & ``Al usar los talonarios de papel hay ocasiones en que la recepcionista coloca a un mismo médico para atender 2 pacientes al mismo horario. Necesitamos que el sistema impida solapar citas.''\\
GER-02 (Gerente) & ``Los cobros los anotamos en un talonario de recibos impresos y al final del día hago el arqueo a mano. Además, los pacientes se confunden o no asisten a la hora adecuada porque no hay un control riguroso de seguimiento.''\\
MED-03 (Médico Tratante) & ``Nosotros pasamos informe al final del día de lo que gastamos en vendajes o gasas para curaciones. A veces se compran insumos de manera tardía y nos quedamos sin material para trabajar.''\\
ADM-04 (Administrador) & ``Guardo las facturas de compras físicas en una carpeta y actualizo el Excel de inventario solo cuando tengo tiempo, lo que genera retrasos y descontrol en la información de existencias.''\\
\bottomrule
\end{longtable}

\section{Procesos del Negocio Observados}

Describa brevemente cómo funciona el flujo actual (AS-IS) o cómo se desea el nuevo flujo simplificado (TO-BE):

\begin{enumerate}[leftmargin=0pt,label={},itemindent=0pt]
    \item[\textbf{Paso 1:}] La recepcionista atiende al paciente presencialmente o por teléfono y anota la cita a mano en un talonario físico de papel.
    \item[\textbf{Paso 2:}] El paciente asiste a su consulta o procedimiento (medicina general, curación de heridas o vendajes) y la recepcionista cobra en efectivo, emitiendo un recibo en papel.
    \item[\textbf{Paso 3:}] El médico utiliza insumos de su consultorio y entrega un informe verbal o manuscrito al final de la jornada con lo consumido.
    \item[\textbf{Paso 4:}] El administrador recibe las facturas físicas de compra de materiales, las guarda en una carpeta y actualiza un archivo de Excel de forma extemporánea cuando dispone de tiempo.
    \item[\textbf{Paso 5:}] Al final de la jornada, la recepcionista realiza un arqueo manual comparando el efectivo recaudado contra las copias del talonario de recibos.
\end{enumerate}

\section{Glosario Inicial del Cliente}

Definir términos clave usados por el cliente para evitar malentendidos (ej.: mercadería en tránsito, merma, stock de seguridad).

\begin{description}
    \item[Solapamiento de horarios:] Conflicto que ocurre cuando se asignan dos o más pacientes al mismo médico dentro del mismo intervalo de tiempo.
    \item[Talonario de recibos:] Libreta de papel físico autocopiativo utilizada en recepción para anotar el cobro de la consulta y la venta de medicamentos.
    \item[Insumos / materiales médicos:] Artículos consumibles utilizados durante la atención, como gasas, vendajes elásticos, apósitos, guantes y productos medicinales.
    \item[Arqueo manual:] Conteo físico del dinero cobrado en recepción, comparado manualmente contra las hojas del talonario de recibos al final del día.
    \item[Compra bajo pedido:] Proceso actual en el que el reabastecimiento de insumos se realiza únicamente cuando el médico notifica que el producto ya se ha agotado.
\end{description}
